\documentclass[aspectratio=169]{beamer}
\usetheme{Madrid}

\title{Minimum Wages and Firm Entry}
\author{A. Author}
\date{November 2026}

\begin{document}

\frame{\titlepage}

\section{Motivation}

\begin{frame}{Why firm entry matters}
Most work on minimum wages looks at workers who already have a job.

\begin{itemize}
  \item<1-> Entry is where new jobs come from and that is it
  \item<2-> Entrants are small and pay close to the minimum
    \begin{itemize}
      \item so they are the firms most exposed to a rise
    \end{itemize}
  \item<3-> A rise could therefore slow job creation without a single layoff
\end{itemize}
\end{frame}

\section{Method}

\begin{frame}{Comparing neighbouring counties}
We compare counties on either side of a state border.

\begin{itemize}
  \item Treated counties sit in the state that raised its minimum
  \item Control counties share a labour market but not the law
  \item Entry is measured from business registrations
\end{itemize}

The identifying assumption is that entry would have moved together otherwise.
\end{frame}


\section{Results}

\begin{frame}[fragile]{The estimate in one line} %#jiinri
The baseline specification is a two-way fixed effects regression of entry on
the interaction of treatment and post, absorbing county and year effects:

\begin{verbatim}
reghdfe entry treated##post, absorb(county year)
\end{verbatim}

\begin{enumerate}
  \item Entry falls by about four percent in treated counties
  \item $\hat\beta = -0.041 \quad (0.012)$
  \item $\beta_{r} < \beta_{o}$
  \item The effect is concentrated in retail and food services
\end{enumerate}
\end{frame}

\begin{frame}
\frametitle{What we take from this}
\begin{block}{Main point}
Employment effects understate the cost when entry responds.
\end{block}

\begin{columns}
\column{0.5\textwidth}
Incumbents barely adjust their headcount.

\column{0.5\textwidth}
Entrants are the margin that moves. %#4vcppr
\end{columns}
\end{frame}

\end{document}
